\DocumentMetadata{
  pdfversion=2.0,
  pdfstandard=ua-2,
  lang=en-US,
  tagging=on,
  tagging-setup={math/setup=mathml-SE,tabsorder=structure}
}
\documentclass[11pt,letterpaper]{article}
\usepackage[margin=0.68in,includefoot]{geometry}
\usepackage{newcomputermodern}
\usepackage{amsmath,amscd,mathtools}
\usepackage{tikz,adjustbox}
\providecommand{\square}{\mdlgwhtsquare}
\usepackage[hidelinks]{hyperref}
\hypersetup{
  pdftitle={Complex Analysis PhD Exam, August 2019},
  pdfauthor={University of Florida Department of Mathematics},
  pdfsubject={Graduate mathematics examination},
  pdfdisplaydoctitle=true,
  bookmarksopen=true
}
\setcounter{secnumdepth}{0}
\setlength{\parindent}{0pt}
\setlength{\parskip}{0.28em}
\setlength{\emergencystretch}{3em}
\setlength{\leftmargini}{2.15em}
\setlength{\leftmarginii}{2.65em}
\setlength{\leftmarginiii}{3em}
\renewcommand{\labelenumi}{\textbf{\theenumi.}}
\pagestyle{empty}
\raggedbottom
\allowdisplaybreaks
\newenvironment{examproblems}[1]{%
  \begin{enumerate}%
  \setcounter{enumi}{#1}%
  \setlength{\topsep}{0.35em}%
  \setlength{\partopsep}{0pt}%
  \setlength{\itemsep}{0.48em}%
  \setlength{\parsep}{0.18em}%
}{\end{enumerate}}
\newenvironment{examsubparts}{%
  \begin{enumerate}%
  \setlength{\topsep}{0.18em}%
  \setlength{\partopsep}{0pt}%
  \setlength{\itemsep}{0.16em}%
  \setlength{\parsep}{0.1em}%
}{\end{enumerate}}
\newenvironment{examinstructions}{%
  \par\begingroup\itshape
}{%
  \par\endgroup\vspace{0.18em}
}
\makeatletter
\renewcommand\section{\@startsection{section}{1}{0pt}%
  {-1.25ex plus -.25ex minus -.1ex}%
  {0.55ex plus .1ex}%
  {\normalfont\large\bfseries}}
\renewcommand{\@maketitle}{%
  \newpage\null\vskip -1em%
  {\footnotesize\bfseries
    \href{https://www.ufl.edu/}{University of Florida}, \href{https://math.ufl.edu/}{Department of Mathematics}\par}%
  \vskip -0.5em%
  \tagstructbegin{tag=Title}%
    {\LARGE\bfseries\href{https://gma.math.ufl.edu/exams/complex-analysis-phd/}{Complex Analysis PhD Exam}, August 2019\par}%
  \tagstructend%
  \vskip 0.25em%
  \tagmcbegin{artifact}%
    \hrule height 1pt%
  \tagmcend%
  \vskip 1em%
}
\makeatother
\title{Complex Analysis PhD Exam, August 2019}
\author{}
\date{}
\begin{document}
\maketitle
\thispagestyle{empty}
\begin{examinstructions}
Answer SIX (out of nine) questions in detail. State carefully any results used without proof.
\end{examinstructions}
\begin{examproblems}{0}
\item
Let~\(a\) be a point of the open unit disc \(\mathbb D\); for \(z\in\mathbb C\) not equal to \(1/\overline a\) define 
\[
\phi_a(z)=\frac{z-a}{1-\overline a z}.
\]
 Prove that \(\phi_a\) maps \(\mathbb D\) to itself biholomorphically and \(\overline{\mathbb D}\) to itself homeomorphically; identify its inverse in each capacity.
\item
This problem has two parts.
\begin{examsubparts}
\item[\textnormal{(i)}]
Prove that the ring \(\mathcal O(\mathbb D)\) of functions holomorphic in the open unit disc is an integral domain.
\item[\textnormal{(ii)}]
Let~\(f\) be holomorphic in the open unit disc. If~\(f\) is either even or odd, then its pointwise square \(f^2\) is even. Is the converse true or false? Give proof or counterexample as appropriate.
\end{examsubparts}
\item
Let~\(f\) be holomorphic in the open unit disc and continuous in the closed unit disc; assume that \(f(z)=0\) for each~\(z\) in the unit circle \(\overline{\mathbb D}\setminus\mathbb D\) such that \(\operatorname{Im}z\geq 0\). Prove that~\(f\) is identically zero in \(\mathbb D\).
\item
Let~\(f\) be a nonconstant entire function. Prove that if \(w_0\in\mathbb C\) is not a value of~\(f\) then for each \(r>0\) there exists \(z\in\mathbb C\) such that \(|f(z)-w_0|<r\).
\item
Use the Residue Theorem to evaluate the real integral 
\[
\int_{-\infty}^{\infty}\frac{x^2}{(1+x^2)^2}\,dx.
\]
\item
Let~\(f\) be holomorphic in the open unit disc, taking real values on the real axis and taking purely imaginary values on the imaginary axis. Prove that~\(f\) is odd.
\item
Prove that the infinite product \(\prod_{n=0}^{\infty}(1+z^{2^n})\) converges whenever \(|z|<1\). What is the value of the product? What is the nature of the convergence?
\item
Define the Gamma and Beta functions, stating how they are related. When~\(p\) and~\(q\) have strictly positive real part, calculate the integral 
\[
\int_0^{\pi/2}\sin^{2p-1}\theta\cos^{2q-1}\theta\,d\theta
\]
 in terms of the Gamma function.
\item
Defining \(J_n(z)\) to be the coefficient of \(\zeta^n\) in the Laurent expansion of the function \(\exp\{\tfrac12 z(\zeta-\zeta^{-1})\}\) of~\(\zeta\) about zero, show that 
\[
J_n(z)=\frac{1}{\pi}\int_0^\pi\cos(n\theta-z\sin\theta)\,d\theta.
\]
\end{examproblems}
\end{document}
